\documentclass[12pt]{article}
\usepackage[utf8]{inputenc}
\usepackage[T1]{fontenc}
\usepackage{hyperref}
\title{Tiv Language Resources for K-12 Curriculum Development}
\author{Compiled by Claude Sonnet 5 (claude-sonnet-5) for Andr\'as Kornai}
\date{}
\begin{document}
\maketitle
\section{Introduction}
Tiv (ISO 639-3 tiv; Glottocode tivv1240) is a Benue-Congo language of Nigeria and Cameroon, written in Latin script. Speaker estimates vary widely (over 4.6 million per Wikipedia; about 15 million including diaspora per the Living Dictionary). My searches found reference works and a modern online dictionary but no primers or textbooks.
\section{Primary Educational Resources}
None found in search results.
\section{Linguistic and Reference Materials}
R.\,C. Abraham, \emph{A Dictionary of the Tiv Language} (1940): \url{https://glottolog.org/resource/reference/id/7621}, \url{https://library.britishmuseum.org/bib/4314}. Orthography and history overview: \url{https://omniglot.com/writing/tiv.htm}. The first Tiv grammar is reported to date from 1933 (no URL seen for it itself).
\section{Religious and Community Resources}
Joshua Project profile (notes continuing need for Christian resources): \url{https://www.joshuaproject.net/index.php/people_groups/15505/NI}. Omniglot reports partial Bible translations; the status of a full Bible was not confirmed.
\section{Digital and Audio Resources}
Tiv Living Dictionary, based on the 2012 revised Tiv orthography of the Tiv Language Studies and Development Association: \url{https://livingdictionaries.app/tiv} and \url{https://livingdictionaries.app/tiv/about}.
\section{Curriculum Development Context}
The reference description says Tiv is taught as a subject in education; I found no confirming policy or syllabus document. A curriculum could build on the 2012 orthography and the Living Dictionary; textbooks and syllabi should be sought from Nigerian education authorities.
\end{document}